% Draft, 2026-11-12. IEEE conference template (IEEEtran, conference mode).
% NOT COMPILED on the authoring machine (no LaTeX installed): build it on Overleaf or with
%   pdflatex main && bibtex main && pdflatex main && pdflatex main
% Sections I-V and VIII are written in full. Sections VI-VII are structured placeholders:
% every \todo{...} needs real data or a decision. Figures and tables are pulled from
% paper/results/, which is a copy of the output folder of
%   python -m analysis.pipeline --data <SD copy> --out paper/results
% Until that folder holds REAL results, the placeholders print instead.
% No number from synthetic data may be moved into Sections VI-VII.
\documentclass[conference]{IEEEtran}
\usepackage[utf8]{inputenc}
\usepackage[T1]{fontenc}
\usepackage{amsmath}
\usepackage{graphicx}
\usepackage{booktabs}
\usepackage{xcolor}
\usepackage{url}
\usepackage{cite}
\DeclareUnicodeCharacter{00B0}{\ensuremath{^\circ}}
\DeclareUnicodeCharacter{00B1}{\ensuremath{\pm}}
\DeclareUnicodeCharacter{00B2}{\ensuremath{^2}}
\DeclareUnicodeCharacter{03B1}{\ensuremath{\alpha}}
\DeclareUnicodeCharacter{03B2}{\ensuremath{\beta}}
\DeclareUnicodeCharacter{03B3}{\ensuremath{\gamma}}

\newcommand{\degC}{\ensuremath{^\circ}C}
\newcommand{\todo}[1]{\textcolor{red}{[TODO: #1]}}
\newcommand{\resultdir}{results}
% A figure from the analysis output, or a labelled box while it does not exist.
\newcommand{\resultfig}[4]{% file stem, width, caption, label
  \begin{figure}[t]\centering
  \IfFileExists{\resultdir/#1.pdf}{\includegraphics[width=#2]{\resultdir/#1.pdf}}%
  {\fbox{\parbox{0.9\columnwidth}{\centering\vspace{6mm}\texttt{\detokenize{#1}.pdf}\\
   produced by \texttt{analysis/report.py} from real data\vspace{6mm}}}}
  \caption{#3}\label{#4}\end{figure}}
% A table written by analysis/report.py (it carries its own caption and label tab:<stem>).
\newcommand{\resulttable}[1]{%
  \IfFileExists{\resultdir/#1.tex}{\input{\resultdir/#1.tex}}%
  {\begin{table}[t]\centering\caption{Placeholder for \texttt{\detokenize{#1}.tex}.}
   \label{tab:#1}\fbox{\parbox{0.9\columnwidth}{\centering
   written by \texttt{analysis/report.py} from real data}}\end{table}}}

\begin{document}

\title{Thermal Spikes on a PV Module in Makkah: Field Characterisation and Real-Time
Detection on a Microcontroller\\
{\large \todo{working title; the supervisor chooses the final one}}}

% Author order as given by Asim on 2026-11-12; to be confirmed with the supervisor.
\author{\IEEEauthorblockN{Asim Alharbi, Abdullah Dahlan, Faisal Al-Fahmi, Ibrahim Al-Khudari,\\
Faisal Al-Maliki, and Abdulilah Al-Haimani}
\IEEEauthorblockA{Department of Mechanical Engineering (Mechatronics)\\
Umm Al-Qura University, Makkah, Saudi Arabia\\
\todo{e-mail addresses}}}

\maketitle

\begin{abstract}
The temperature of a photovoltaic (PV) module does not only drift with the day; it also
rises in steps of several degrees within minutes, for example when the sun returns after
a cloud. Such fast rises are well studied as a modelling problem, but they are rarely
described as events: how often they occur at a site, how large they are, what drives
them and what they cost in power. This paper characterises these \emph{thermal spikes},
defined as a rise of the module temperature of at least 3\,\degC{} within 10\,min, on a
100\,W module in Makkah, Saudi Arabia, using a low-cost instrument built around an ESP32
microcontroller: five back-sheet temperature sensors, an in-plane irradiance sensor and a
two-branch switched resistive load that both holds the module at its maximum-power point
and sweeps its current--voltage curve. The temperature coefficient of power, absent from
the module's documentation, is measured outdoors from these sweeps. Three causal detectors that use only 10-s samples available on the microcontroller
are scored event by event: a temperature-rise
threshold, a CUSUM on the temperature increment, and a small gradient-boosted tree model.
The comparison asks how early each detector raises its alarm and how many false alarms
that costs. \todo{Two to three sentences of results from Section VI: observation period,
spikes per day by cause, power cost per spike, and the lead time and false-alarm rate of
each detector.}
\end{abstract}

\begin{IEEEkeywords}
photovoltaic module temperature, thermal transient, event detection, CUSUM, gradient
boosted trees, TinyML, microcontroller, temperature coefficient, I--V curve
\end{IEEEkeywords}

\section{Introduction}
The power of a crystalline-silicon PV module falls as its temperature rises; reported
field values of the power temperature coefficient lie between 0.30 and 0.45\,\%/\degC{}
in absolute value \cite{kaldellis2014}, and the many correlations in use are reviewed in
\cite{skoplaki2009}. In a hot, dry climate the module spends most of the day far above
the 25\,\degC{} of its rating, so its temperature is a first-order variable of its yield.

Module temperature is usually discussed as a level: models give it from irradiance, air
temperature and wind at steady state \cite{king2004,faiman2008}, and cooling systems
switch on when it crosses a set value \cite{moharram2013}. But the temperature also
moves fast. A module has a thermal time constant of about six minutes \cite{herteleer2023},
so a step in irradiance or a lull in the wind appears as a rise of several degrees
within minutes. A growing body of work models these dynamics so that predicted
temperature follows measured temperature at second-to-minute scales
\cite{jones2001,armstrong2010,lobera2013,veldhuis2015,barry2020,prilliman2020,herteleer2023}.

What that work does not provide, in the literature we read, is a description of the
fast rises as events. For a given module and site we could not find how often they
occur, how large they are, which share follows an irradiance step, and what they cost in
power; nor could we find such measurements for western Saudi Arabia. The question
matters for any action that should be taken \emph{during} a rise, such as a short burst
of cooling: the action is only useful if the rise is recognised while most of it is
still to come.

This leads to the question of this paper: \emph{how often do thermal spikes occur on a
fixed 100\,W module in Makkah during the measurement period, what drives them, what do
they cost in power, and how early can a microcontroller detect them?} We call a thermal
spike a rise of the module temperature of at least 3\,\degC{} within 10\,min, whatever
its cause; the rise that follows the return of the sun after a cloud is included.

The contributions are deliberately modest.
\begin{enumerate}
\item A documented low-cost instrument: an ESP32 logger with five back-sheet sensors,
an in-plane silicon irradiance sensor and a two-branch switched resistive load that
holds the maximum-power point and sweeps the I--V curve, with an uncertainty budget and
a comparison against a commercial data logger.
\item A characterisation of thermal spikes on one module over a stated period: frequency,
size, duration, cause (irradiance-driven or not) and power cost, the last through a
temperature coefficient measured on this module by an outdoor transient method
\cite{king1997}.
\item A comparison of three causal detectors, built for the microcontroller and scored at
event level \todo{say ``running on the microcontroller'' only once all three do; see
Section~\ref{sec:ondevice}}.
The comparison is framed around \emph{detection lead time and false alarms}, not around
a headline accuracy: with a rise defined by a threshold, a threshold rule will find
nearly every event sooner or later, so the useful question is how much earlier another
detector fires and at what cost in false alarms.
\end{enumerate}
To the authors' knowledge, fast temperature rises of a PV module have not been reported
as discrete events with their frequency, cause and power cost from field measurements in
Makkah, and detectors of such events have not been compared in real time on a low-cost
microcontroller. This statement rests on a limited search, described in Section~\ref{sec:related}.

The paper does not claim a climate statistic (one module, one site, one period), a
cooling benefit, a new thermal model or a new I--V tracer topology, and it reports the
detector comparison whichever way it comes out.

\section{Related Work}\label{sec:related}
\subsection{Module temperature: steady state and dynamics}
The Sandia model gives the back-surface temperature as
$T_m = E\,\exp(a+b\,W\!S)+T_a$ and the cell temperature as $T_c = T_m + (E/E_0)\,\Delta T$,
with $\Delta T = 3$\,\degC{} for an open-rack glass/cell/polymer-sheet module and a
stated accuracy of about $\pm5$\,\degC{} \cite{king2004}. Faiman showed on seven module
types that a modified Hottel--Whillier--Bliss equation predicts module temperature to
within the cell-to-cell differences inside a module \cite{faiman2008}; its
implementation in pvlib \cite{pvlib2023} is the thermal model used in this work. Wind is
decisive: at a test site in Bolzano, models that include wind predicted cell temperature
better than those that do not for most technologies \cite{schwingshackl2013}.

Steady-state models fail when conditions change faster than the module follows. Jones
and Underwood treated the module as a non-steady energy balance and reported agreement
within 6\,K for 95\,\% of the time under cloudy conditions \cite{jones2001}. Armstrong
and Hurley argued that a steady-state model cannot be justified during rapidly changing
conditions and proposed an RC model \cite{armstrong2010}; Torres Lobera and Valkealahti
validated a dynamic model on three months of 1-s data \cite{lobera2013}. Empirical and
filter-based corrections followed: a model for transient weather at 1-min resolution
\cite{veldhuis2015}, an exponential smoothing kernel that halved the error of a
steady-state model \cite{barry2020}, a weighted moving average over the previous 20\,min
\cite{prilliman2020}, and a simplified energy balance checked on a case with passing
clouds \cite{hannouch2024}. Herteleer et al.\ analysed 24 datasets from 15 sites and
report a thermal time constant of $(6.3\pm1.0)$\,min for fixed-mount systems
\cite{herteleer2023}; we use this value for the detection window and for the cause label.

All of these treat a fast change as something to model. None of those we read reports
the rises as events with their frequency, size, cause and power cost.

\subsection{Temperature coefficients measured outdoors}
King et al.\ describe an outdoor procedure: the module is shaded until it is near air
temperature, then uncovered, and I--V curves are measured while it heats, with the
module temperature taken as the average of several sensors on the rear surface. They
name two difficulties: non-uniform temperature during the transient, and a back-surface
temperature that does not represent the cell temperature \cite{king1997}. Coefficients
have also been estimated from long-term field data with irradiance and spectral
corrections \cite{paudyal2021}. IEC~60891 defines the correction procedures and the
determination of their factors \cite{iec60891}; we describe our own procedure and do not
claim conformity with the standard.

\subsection{Low-cost monitoring and I--V tracing}
Methods to obtain I--V curves are reviewed in \cite{duran2008,zhu2020,morales2021}; the
last lists six topologies, one of which uses the duty cycle of a DC--DC converter to
emulate a resistor. Low-cost tracers have been built on open microcontroller platforms
\cite{papageorgas2015} and with a relay-switched resistor bank \cite{machado2025}.
Low-cost monitoring systems have been compared with commercial instruments
\cite{paredes2018} and designed to the accuracy requirements of the monitoring standard
\cite{lopezvargas2019,iec61724}. Our load switches resistors by pulse-width modulation,
without an inductor, and uses the same circuit to hold the maximum-power point and to
sweep the curve. We did not find this arrangement in the papers we read; we do not
claim it as new.

\subsection{Event detection and learning on microcontrollers}
Sequential change detection goes back to Page's CUSUM \cite{page1954}; methods for time
series are surveyed in \cite{aminikhanghahi2017}. Benchmarks for real-time detectors
score how early an anomaly is found \cite{lavin2015}, and precision and recall have been
extended from points to ranges \cite{tatbul2018}; our metrics belong to this family.
Gradient-boosted trees \cite{ke2017} remain strong on medium-sized tabular data
\cite{grinsztajn2022}, and tree models can run in a few kilobytes of memory
\cite{kumar2017}; the field is reviewed in \cite{ray2022}.

In PV, the closest systems we read differ in kind. Novak and Landkamer use ESP32 nodes
to sense and transmit, while the decision to cool is taken on a single-board computer
from the panel temperature, the coolant state and weather trends \cite{novak2026}.
Mellit classifies infrared images of modules with a convolutional network on a
single-board computer \cite{mellit2022}. We did not find a detector of thermal
transients that runs on the microcontroller itself and is scored at event level against
simple baselines. \todo{Before submission: read \cite{herteleer2023}, \cite{king1997}
and \cite{novak2026} in full and repeat the search for on-microcontroller detection;
most works above were read at abstract level.}

\section{System}
\subsection{Module and site}
The module is a 100\,W mono-crystalline unit of the 12\,V class. Its seller lists
$V_{mp}=18$\,V, $I_{mp}=5.55$\,A, $V_{oc}=21.5$\,V and $I_{sc}=5.88$\,A; no temperature
coefficients, operating-cell temperature or cell count are published, which is why the
coefficient is measured here. It is mounted on a fixed open rack at a tilt of
20$^\circ$ facing true south \todo{tilt and azimuth as measured at installation;
coordinates and a photograph of the site}.

\subsection{Sensors and logger}
Table~\ref{tab:instrument} lists the measured quantities. An ESP32 reads all sensors
every 10\,s and writes one row to a microSD card with a UTC timestamp from a
temperature-compensated real-time clock. Five digital temperature sensors are fixed to
the back sheet behind cell centres; their mean is the module temperature $T_m$. The
irradiance sensor is mounted in the plane of the module. A button on the logger writes a
time mark for every deliberate action at the rig; a daily log sheet gives each mark its
meaning. \todo{State which sensors were calibrated in the common water bath and against
which reference.}

\begin{table}[t]
\caption{Measured quantities and instruments.}
\label{tab:instrument}
\centering
\begin{tabular}{lll}
\toprule
Quantity & Instrument & Interface \\
\midrule
Back-sheet temperature ($\times5$) & DS18B20 & 1-Wire \\
Air temperature, humidity & BME280, shielded & I$^2$C \\
In-plane irradiance & PYR20 silicon sensor & RS-485 Modbus \\
Module voltage and current & INA226, 5\,m$\Omega$ shunt & I$^2$C \\
Time & DS3231 clock, UTC & I$^2$C \\
\bottomrule
\end{tabular}
\end{table}

\subsection{Measurement load}
The module must work at its maximum-power point (MPP) for its temperature and its power
to be those of a module in service, and its I--V curve must be swept to measure the
maximum power directly. One circuit does both. Two identical branches, each a
1.5\,$\Omega$ resistor bank switched by its own MOSFET at 20\,kHz, are connected across
the module behind an input capacitor. With duty cycle $D$ the module sees an effective
resistance
\begin{equation}
R_\mathrm{eff} = R/D ,
\end{equation}
with $R=1.5\,\Omega$ when branch~A switches alone and $R=0.75\,\Omega$ when both switch
together. Branch~A alone holds the module at the voltage of the last measured MPP.
Once a minute both branches sweep the duty cycle for about 1.6\,s: a coarse pass from
open circuit down to about 4.4\,V, then a fine pass around the maximum. The current
sensor sits between the module and the capacitor, so each recorded pair $(V,I)$ lies on
the module's own curve. Using the second branch only during sweeps halves the ripple
current the capacitor carries all day. \todo{Final part values, the measured settling
time and the bench results (docs/paper/06); a one-line schematic figure.}

\subsection{On-device detection}\label{sec:ondevice}
The detectors of Section~\ref{sec:detectors} use only the current and earlier 10-s samples.
The tree model and its features run on the logger's microcontroller; the model is exported
to plain C \cite{m2cgen}. \todo{The threshold and CUSUM detectors are computed offline from
the same rows in the present firmware. Port both to the logger (a few lines of C each) and
log their alarms, or state here that only the tree model ran on the device.} \todo{Time per update, flash and RAM
measured on the board with the final model.}

\section{Methodology}
\subsection{Definition of a thermal spike}
Reference labels are made offline, with the whole day in view; the detectors see only
past samples. This separation matters: because a spike is defined by a threshold on the
rise, a causal threshold rule fires at the moment the rise \emph{reaches} the threshold
and never before. The comparison is therefore about how early and how reliably a
detector fires while the rise is still developing.

Let $T_s$ be the module temperature smoothed with a centred 3-sample median and a
centred 1-min mean. With a window $W=10$\,min,
\begin{equation}
R(t) = T_s(t) - \min_{t-W\le u\le t} T_s(u).
\end{equation}
A spike exists when $R(t)\ge\theta$, $\theta=3$\,\degC. Its onset $t_\mathrm{on}$ is the
latest minimum of $T_s$ in the window that ends at the first crossing $t_\theta$; its
peak $t_\mathrm{pk}$ is the first maximum after $t_\theta$ that is not exceeded during
the next 5\,min; its size is $A=T_s(t_\mathrm{pk})-T_s(t_\mathrm{on})$. The threshold is
three times the normal heating rate of about 0.1\,\degC/min reported in
\cite{moharram2013}, and the window is about 1.6 thermal time constants
\cite{herteleer2023}. Rises between 2 and 3\,\degC{} are recorded as minor rises and are
neutral in scoring. Events are kept only when the sun is more than about 11$^\circ$
above the horizon and outside excluded periods: 30\,min after each logged cleaning, and
the shade-and-uncover runs of Section~\ref{sec:gamma}.

\subsection{Cause and origin}
Wind is not measured, so the cause is labelled only as irradiance-driven or not. Let
$\Delta G$ be the difference between the mean irradiance during the event and the lowest
1-min mean in the 10\,min before it. A first-order lag gives the rise this step explains,
\begin{equation}
\Delta T_G = \frac{\Delta G}{U_\mathrm{eff}}
\left(1-e^{-(t_\mathrm{pk}-t_\mathrm{on})/\tau}\right),
\end{equation}
with $\tau=6.3$\,min \cite{herteleer2023} and $U_\mathrm{eff}$ fitted on event-free
clear-sky data of the campaign \todo{fitted value; starting value 35\,W/(m$^2$K)}. The
event is labelled \emph{irradiance} when $f_G=\Delta T_G/A\ge 2/3$,
\emph{non-irradiance} when $f_G\le 1/3$, and \emph{mixed} otherwise. A second field
records the origin: natural, or one of the controlled tests (a shade mesh of about 20\,\%
over module and sensor, removed after 20\,min; a windbreak for 20\,min). Controlled
tests have a known cause, so they also serve to check the cause rule.

\subsection{Temperature coefficient and power cost}\label{sec:gamma}
The coefficient $\gamma$ of the maximum power is measured by the outdoor transient
method of \cite{king1997}. Near solar noon on a clear day the module, but not the
irradiance sensor, is shaded for 25\,min and then uncovered; the load sweeps the I--V
curve every 15\,s for 20\,min while the module heats. For each sweep the maximum power
$P_{mp}$ is the vertex of a parabola fitted to the points within 3\,\% of the largest
power. The cell temperature is $T_c=T_m+\Delta T\,G/G_0$ with $\Delta T=3$\,\degC{}
\cite{king2004}; the first 2\,min after uncovering are dropped, because the cells lead
the back sheet there. Sweeps during which the irradiance changed by more than 2\,\% are
rejected. A least-squares line through
$P^*=P_{mp}\,G_\mathrm{ref}/G$ against $T_c-25$ gives intercept $a$ and slope $b$, and
\begin{equation}
\gamma = b/a .
\end{equation}
As a check that does not use the irradiance sensor, the same regression on
$P_{mp}/I_{sc}$ gives $\gamma-\alpha$, where $\alpha$ is the coefficient of the
short-circuit current. The power cost of a spike is then
\begin{equation}
\Delta P/P = \gamma\,\Delta T_c,\qquad
\Delta T_c = A + \Delta T\,\Delta G/G_0 .
\end{equation}
It is computed, not read from the measured power, for the reason given in Section~\ref{sec:uncertainty}.

\subsection{Detectors}\label{sec:detectors}
Three detectors use the causal 3-sample median of $T_m$; a fourth, deliberately trivial
one looks at irradiance only. All are silent for 10\,min after an alarm.
\begin{itemize}
\item \emph{Threshold}: alarm when the rise over the last 10\,min is at least
$\theta_a$. No irradiance condition is applied, since irradiance-driven rises count.
\item \emph{CUSUM}: $S_k=\max(0,\,S_{k-1}+(T_k-T_{k-1})-\kappa)$ on the temperature
increment, alarm when $S_k\ge h$ \cite{page1954}.
\item \emph{Gradient-boosted trees}: at most 50 trees of depth four \cite{ke2017} on
ten features that the microcontroller computes: temperature rise over 2, 5 and 10\,min;
relative irradiance change over 2\,min and range over 10\,min; clear-sky index;
cosine of the solar zenith angle; temperature excess over air per unit irradiance; and
the residual of a first-order thermal model \cite{faiman2008} with its CUSUM. The
target is the first 10\,min of every spike. The alarm requires the output to exceed a
threshold on two consecutive samples.
\item \emph{Irradiance step}: alarm when the irradiance has risen by a fraction $x$ or
more within 2\,min. It never sees the temperature. It is included as a check on the
tree model: for spikes that follow a returning sun, a one-line rule on the irradiance may
be all that is needed, and the tree model must be shown to add something beyond it.
\end{itemize}

\subsection{Evaluation}
Days are split chronologically into training (60\,\%), validation (20\,\%) and test
(20\,\%), by whole local days. Every setting ($\theta_a$; $\kappa$, $h$; the output
threshold; $x$) is chosen on the validation days and the test days are scored once.

An alarm matches a spike when it falls between the onset and 2\,min after the peak. The
first match is a true positive; a spike without a match is missed; an alarm outside
every spike, minor rise and excluded period is a false alarm. We report precision,
recall and $F_1$, and recall and $F_1$ per cause (with all false alarms in the
precision, since a false alarm has no cause). The quantities the comparison rests on
are:
\begin{itemize}
\item the \emph{detection delay} $t_\mathrm{alarm}-t_\mathrm{on}$;
\item the \emph{lead} $t_\theta-t_\mathrm{alarm}$ over the moment the rise reaches
3\,\degC, positive when the alarm comes first;
\item the \emph{share of the rise remaining} at the alarm,
$(T_s(t_\mathrm{pk})-T_s(t_\mathrm{alarm}))/A$;
\item \emph{false alarms per day}, counted per 12 valid daytime hours: excluded periods
and samples the detector did not receive are removed from the observation time.
\end{itemize}
A detector that alarms more readily also alarms earlier, so delays are compared in two
ways. First at the settings with the best validation $F_1$. Second at equal
false-alarm budgets: for each budget, each detector is given the setting with the
shortest median delay on the validation days among those that stay within the budget
and still find at least 90\,\% of the validation spikes; that setting is then scored
on the test days. We also show, for every setting, delay and recall against false alarms
per day.
Intervals are 95\,\% percentile intervals from 2000 bootstrap resamples of whole days;
differences between detectors are computed on the same resampled days. A difference
whose interval contains zero is reported as not detectable. The alarms written by the
logger are compared with the same rows replayed offline.

\subsection{Verification of the analysis chain on synthetic data}
Before any field data existed, the complete chain was run on 150 synthetic days written
in the logger's file formats. This verifies the software only. The synthetic module
follows the same first-order thermal model that one of the features uses, its spikes are
large (median 6.6\,\degC) and most follow a deep, abrupt cloud, and the rate of events
is set by the generator, not by any climate. On the 30 synthetic test days the three
temperature-based detectors found nearly every spike and their $F_1$ values (0.97 to
0.99) could not be told apart. At the settings with the best validation $F_1$ the tree
model alarmed a median of 1.8\,min after onset, against 7.3\,min for the threshold rule
and 5.7\,min for CUSUM. Two checks qualified this. At equal false-alarm budgets the lead
over the threshold rule remained (4 to 5\,min), but a CUSUM set for speed alarmed after
2.0\,min and its difference from the tree model was not detectable; it then found
87\,\% of the spikes against 95 to 100\,\%. And a rule that alarms when the irradiance
rises by half within 2\,min was as early as the tree model on irradiance-driven spikes
(0.5 against 0.7\,min); it found 87\,\% of those and none of the others, with 0.39 false
alarms per day against none. These numbers are not results of this paper. They set what
the field comparison has to show: not that a learned model is earlier, but whether it
adds recall and removes false alarms beyond what a fast CUSUM and an irradiance rule
give.
\todo{Decide with the supervisor whether this subsection stays; if space is short, keep
its first two sentences only.}

\section{Measurement Uncertainty}\label{sec:uncertainty}
Uncertainties follow the GUM \cite{gum2008}: datasheet limits are treated as rectangular
distributions and combined in quadrature; Table~\ref{tab:uncertainty} gives the result.
Entries marked with a dagger rest on an assumed value until the measurement named in the
text replaces it.

\begin{table}[t]
\caption{Standard uncertainties ($k=1$) before field calibration.}
\label{tab:uncertainty}
\centering
\begin{tabular}{lll}
\toprule
Quantity & Absolute & Change within 10\,min \\
\midrule
$T_m$ (mean of five) & 0.14--0.31\,\degC & 0.03--0.05\,\degC$^\dagger$ \\
$T_c$ from $T_m$ & adds 0.5\,\degC$^\dagger$ & only if $G$ changes \\
Air temperature & 0.65\,\degC$^\dagger$ & 0.1\,\degC$^\dagger$ \\
Irradiance & 3.4--6.7\,\% & \todo{measured} \\
Voltage & 0.11\,\% & 0.01\,\% \\
Current & 0.59\,\%$^\dagger$ & 0.01\,\% \\
Power & 0.60\,\%$^\dagger$ & 0.02\,\% \\
\bottomrule
\end{tabular}
\end{table}

\emph{Temperature.} Each sensor is specified to $\pm0.5$\,\degC; the mean of five lies
between 0.14\,\degC{} (independent errors) and 0.31\,\degC{} (fully correlated). A rise
within one event is a difference of readings minutes apart, so offsets cancel and only
resolution and noise remain: about 0.05\,\degC, sixty times smaller than the 3\,\degC{}
that defines a spike. The cell-to-back offset \cite{king2004} cancels at steady
irradiance and enters an irradiance-driven rise through $\Delta G$.

\emph{Irradiance.} The sensor is a silicon device with a stated accuracy of 5\,\% and a
directional error of 3\,\%; whether the 5\,\% refers to the reading or to full scale is
not stated \todo{seller's answer}, which gives the range in the table. Irradiance is
therefore used only for steps, ratios and stability checks, never as an absolute
reference.

\emph{Electrical quantities.} The current is limited by the tolerance of the shunt,
assumed to be 1\,\% until it is compared with a reference ammeter. Changes over minutes
are resolved to about 0.02\,\% of power.

\emph{Consequence.} A 3\,\degC{} spike changes the power by roughly one percent, which
the electrical chain resolves easily but which cannot be separated from the effect of
irradiance when irradiance is known to a few percent. The power cost is therefore
computed as $\gamma\,\Delta T_c$, and its uncertainty is that of $\gamma$.

\emph{Temperature coefficient.} Scatter within one run contributes about 2\,\% of
$\gamma$. The systematic terms are larger: the cell-to-back offset builds from zero to
about 2.7\,\degC{} during a run, which would overstate $|\gamma|$ by about 13\,\% over a
20\,\degC{} span if ignored and leaves about 2.6\,\% after correction; a drift of
0.5\,\% in the irradiance reading during a run would bias $\gamma$ by about 6\,\% for a
coefficient of $-0.4$\,\%/\degC. The reported value is the mean over runs on several
days, with uncertainty $\sqrt{s^2/n+u_B^2}$, $s$ being the spread between runs.
\todo{Replace the estimates of this paragraph with the measured spread.}

\emph{Time.} The clock drifts by at most 0.17\,s per day between 0 and 40\,\degC; the
offset is checked weekly. Within one row the readings are taken within about one second.

\section{Results}
\emph{This section holds no results yet. Each item names the figure or table that
\texttt{analysis/pipeline.py} writes and the sentence it must support.}

\subsection{Instrument validation}
\resulttable{table5_validation}
\resultfig{fig5_validation}{\columnwidth}{ESP32 minus reference logger against their
mean, back-sheet sensor T3. \todo{dates}}{fig:validation}
\todo{From Table~\ref{tab:table5_validation}: clock lag; bias and limits of agreement of
temperature, of the 10-min rise, of irradiance (same sensor, two read-outs: this checks
the read-out, not the sensor), voltage and current; which acceptance criteria were met.
State the reference logger's model and its own accuracy.}

\subsection{Thermal spikes}
\resultfig{fig1_example_day}{\columnwidth}{One day: module and air temperature,
irradiance, labelled spikes by cause, and the alarms of the four detectors.
\todo{date}}{fig:day}
\resultfig{fig2_spike_statistics}{\columnwidth}{Spike size by cause, and number of
spikes per day.}{fig:stats}
\resulttable{table1_spike_statistics}
\resulttable{table6_cause_check}
\resulttable{table9_cause_sensitivity}
\todo{Observation period, valid days and hours. From Table~\ref{tab:table1_spike_statistics}:
natural spikes per day with a 95\,\% interval, by cause; size, duration and rate. Minor
rises. From Table~\ref{tab:table6_cause_check}: agreement of the cause rule with the
controlled tests. Natural and controlled events are reported separately. From
Table~\ref{tab:table9_cause_sensitivity}: how many spikes change cause when the two
limits (1/3 and 2/3) are moved, and whether any conclusion depends on them.}

\subsection{Temperature coefficient and power cost}
\resultfig{fig4_gamma}{\columnwidth}{Irradiance-corrected maximum power against cell
temperature for all shade-and-uncover runs.}{fig:gamma}
\resulttable{table4_gamma}
\todo{$\gamma$ with expanded uncertainty, number of runs and days; the
$I_{sc}$-normalised value; $\alpha$ and $\beta$; comparison with the field range of
\cite{kaldellis2014} and the values of \cite{king1997}. Then the power cost per spike
(median and range) and summed per day, from the last column of
Table~\ref{tab:table1_spike_statistics}.}

\subsection{Detection}
\resulttable{table2_detection}
\resulttable{table3_per_cause}
\resultfig{fig3_detection_delay}{\columnwidth}{Share of spikes detected against time
from onset; the dotted line marks the median time at which the rise reaches
3\,\degC.}{fig:delay}
\resulttable{table8_matched_false_alarms}
\resultfig{fig6_delay_vs_false_alarms}{\columnwidth}{Left: shortest median delay each
detector reaches with at most the given number of false alarms per day, among its
settings that find at least 90\,\% of the spikes; open markers are the $F_1$-tuned
settings. Right: highest recall reached within the same number of false alarms. Test
days.}{fig:tradeoff}
\todo{Number of spikes in training, validation and test days; tuned settings. Lead
with the two quantities the comparison rests on: (i) lead time and share of the rise
remaining, from Table~\ref{tab:table2_detection} and Fig.~\ref{fig:delay}, with the
paired difference between the tree model and each baseline and its interval; (ii) false
alarms per day with intervals. Then the same comparison at equal false-alarm budgets
(Table~\ref{tab:table8_matched_false_alarms}, Fig.~\ref{fig:tradeoff}), because a lead
bought with more false alarms is not a lead: say whether the order of the detectors and
the size of the lead survive. Then precision, recall and $F_1$, saying plainly whether
their differences are detectable. Then recall and delay per cause from
Table~\ref{tab:table3_per_cause}. State the result for the irradiance-step rule plainly:
its recall and delay on irradiance-driven spikes against the tree model's, and that it
cannot find spikes without an irradiance step; the tree model's case rests on what it
adds for those.
If fewer than 50 training and 20 test spikes were collected, the threshold and CUSUM are
the main result and the tree model is exploratory.}

\subsection{Agreement on the device}
\resulttable{table7_parity}
\todo{Alarm mismatches between the logger and the offline replay of the same rows;
detector blackouts; time per update and memory on the board. Requires the on-board
replay to be re-run with the final model.}

\section{Discussion}
\todo{Write after the results. Points to cover, each tied to a table or figure:}
\begin{itemize}
\item \todo{What the lead time is worth: with a time constant of about 6\,min, an alarm
$x$ minutes earlier leaves $y$\,\% of the rise still to come
(Table~\ref{tab:table2_detection}). Say for which later use this matters.}
\item \todo{Whether the learned detector earns its place. The irradiance-step rule is
in Tables~\ref{tab:table2_detection}--\ref{tab:table8_matched_false_alarms}: where it is
as early as the tree model (irradiance-driven spikes), say so; the tree model is then
justified only by the spikes the rule cannot see and by its false alarms. If that gain is
not detectable in the field data, recommend the simple rules.}
\item \todo{False alarms: what one false alarm per day would cost a later action.}
\item \todo{Natural against controlled events: do they differ in size and rate?}
\item \todo{How rare natural spikes were, and what that means for the question.}
\end{itemize}
\emph{Limitations}, to be kept whatever the results:
one module, one site and a short period, so no climate statistic; a silicon irradiance
sensor of 5\,\% stated accuracy; back-sheet rather than cell temperature; no wind
measurement, so the cause is only irradiance-driven or not; controlled events are not
natural events; reference labels and the threshold detector share one signal, which the
hindsight-versus-causal design addresses but does not remove; the tree model's internal
thermal model uses literature parameters, not ones fitted to this module
\todo{unless refitted}.

\section{Conclusion}
\todo{Three or four sentences after the results: how often spikes occurred and what
drove them; their power cost; how much earlier than a threshold rule a detector on a
microcontroller raised its alarm, and at how many false alarms per day; what the next
step is.} The instrument, the definitions and the analysis code are documented so that
the measurement can be repeated on another module or site.

\section*{Acknowledgment}
\todo{Supervisor, laboratory access, workshop that built the stand, lender of the data
logger. State the use of AI tools as the venue requires.}

\bibliographystyle{IEEEtran}
\bibliography{refs}

\end{document}
